\section{Soundness and Relative Completeness}
\label{sec:metatheory}

We prove soundness and relative completeness by translating certificates in
both directions, following the approach of Seidel~\cite[Chapter~6]{seidel2026thesis}. The only
control-flow parameter is a bound on the number of edges in one body
execution; one backward recursion handles every internal location.
All sums below range over the positive-mass outcomes of their distributions.

\subsection{One Iteration and Its Support}

Fix $C=\while{b}{B}$ with loop-free $B$. Write $\ell_g$ for its guard,
$\ell_{\rm in}$ for the body entry, and $\ell_{\rm exit}$ for the program
exit. In $\mathcal T[C]$, the body exit is identified with $\ell_g$;
there is no additional back edge. Skip is represented by one identity
edge. Let $L_B\geq1$ bound the number of \emph{edges} in any complete body
path. Every such path is finite, and from a fixed state there are only
finitely many paths. The body therefore preserves mass.

Let $p_*$ be the minimum of $1$ and all positive probabilistic edge
probabilities in $\mathcal T[B]$. Zero-probability edges are discarded.
Because the program has finitely many fixed choices, $p_*>0$, and every
feasible body path has probability at least
\begin{equation}
 \varepsilon_B=p_*^{L_B}>0.
 \label{eq:body-path-bound}
\end{equation}
This bound is independent of the valuation and of the certificates.

\paragraph{Support closure.}
If $I$ is a distributional invariant for $C$ and an initial set of distributions $M$, then $S=\supp(I)$ contains
$\supp(M)$ and is closed under enabled body executions. Indeed, for
$\sigma\in S\cap\llbracket b\rrbracket$, choose $\mu\in I$ with
$\mu(\sigma)>0$. Whenever $\post{B}(\dirac{\sigma})(\sigma')>0$,
linearity gives
\[
 \bigl([\neg b]\mu+\post{B}([b]\mu)\bigr)(\sigma')
 \geq\mu(\sigma)\post{B}(\dirac{\sigma})(\sigma')>0.
\]
The distribution on the left belongs to $I$, so $\sigma'\in S$.
Conversely, if $S$ contains $\supp(M)$ and has this closure property,
then $\{\mu\in\subdists\mid\supp(\mu)\subseteq S\}$ is a
distributional invariant: each positive output mass comes from an input
state in $S$, or from terminal mass retained unchanged.

\begin{lemma}[Contracting an iteration]
\label{lem:contract-v}
\label{lem:contract-u}
Suppose $\mathit{Inv},V,U$ satisfy the pCFG premises in
\Cref{thm:pcfg-rule} for $\mathcal T[C]$. At every enabled guard
configuration $(\ell_g,\sigma)\in\mathit{Inv}$,
\begin{align}
 \sum_{\sigma'}\post{B}(\dirac{\sigma})(\sigma')
                V(\ell_g,\sigma')&\leq V(\ell_g,\sigma),
 \label{eq:contract-v}\\
 \sum_{\sigma':\,U(\ell_g,\sigma')<U(\ell_g,\sigma)}
       \post{B}(\dirac{\sigma})(\sigma')&\geq\varepsilon_B.
 \label{eq:contract-u}
\end{align}
\end{lemma}
\begin{proof}
Backward induction over the acyclic body composes the local expectation
inequalities: a deterministic step substitutes its unique successor,
and a probabilistic step takes their weighted sum. Including the enabled
guard edge and using the operational correspondence from
\Cref{app:translation} gives \eqref{eq:contract-v}.

For $U$, the guard edge strictly decreases its value. At each subsequent
nonterminal configuration choose a positive-probability successor with
smaller $U$: the unique successor at a deterministic location, and a
decreasing successor at a probabilistic location. Such a successor exists
by the pCFG progress premise, applied with $r=V(\ell,\sigma)$ at the
current configuration. Acyclicity brings this path back to the guard in
at most $L_B$ body edges. Its endpoint has smaller $U$ than its starting
guard configuration, and its probability is at least $\varepsilon_B$
by \eqref{eq:body-path-bound}. This proves \eqref{eq:contract-u} without
requiring the path to remain in any fixed sublevel set.
\end{proof}

\subsection{Soundness}

\begin{theorem}[Soundness]
\label{thm:soundness}
The program-level rule in \Cref{thm:program-rule} is sound.
\end{theorem}
\begin{proof}
Assume certificates $I,v,u$. Let $\mathit{Inv}$ be the configurations
reachable in $\mathcal T[C]$ from $(\ell_g,\sigma)$ with
$\sigma\in\supp(M)$. Support closure ensures that every reachable
guard valuation belongs to $\supp(I)$. Set $a=L_B+1$ and prescribe
\[
 \begin{array}{ll}
 V(\ell_g,\sigma)=v(\sigma)+1,
   &U(\ell_g,\sigma)=a\,u(\sigma)+1,\\
 V(\ell_{\rm exit},\sigma)=0,
   &U(\ell_{\rm exit},\sigma)=0.
 \end{array}
\]
At internal body configurations $\gamma$, define both functions by
backward recursion, stopping at the next guard visit:
\begin{equation}
 V(\gamma)=\sum_{\gamma'\in\Succ(\gamma)}
                  P(\gamma,\gamma')V(\gamma'),
 \qquad
 U(\gamma)=1+\min_{\gamma'\in\Succ(\gamma)}U(\gamma').
 \label{eq:lift-recursion}
\end{equation}
Here $P$ is the one-step transition probability, equal to $1$ for a
unique deterministic successor; $\Succ$ contains only positive-probability
successors. These are finite, well-defined recursions because the body
is acyclic. All nonterminal $V$- and $U$-values are at least $1$, and
$U$ is natural-valued.

\paragraph{Local inequalities.}
Inside the body, \eqref{eq:lift-recursion} gives the supermartingale
inequality with equality. It also gives strict variant descent on every
deterministic edge and on a successor of probability at least $p_*$ at
every probabilistic location. At a guard with $\neg b$, both functions
strictly decrease to zero. At a guard with $b$, backward expansion and
mass preservation give
\[
 V(\ell_{\rm in},\sigma)
 =\sum_{\sigma'}\post{B}(\dirac{\sigma})(\sigma')
                         (v(\sigma')+1)
 \leq v(\sigma)+1=V(\ell_g,\sigma).
\]
For $U$, the source decrease premise supplies a feasible body path to
some $\sigma'$ with $u(\sigma')\leq u(\sigma)-1$. Along \emph{any}
body path of $n$ edges, the minimum recursion gives
$U(\gamma_0)\leq U(\gamma_n)+n$. Thus, for this path, $n\leq L_B$
and
\begin{align*}
 U(\ell_{\rm in},\sigma)
 &\leq a\,u(\sigma')+1+n\\
 &\leq a\,(u(\sigma)-1)+1+L_B
  =a\,u(\sigma)<U(\ell_g,\sigma).
\end{align*}
The extra unit in $a=L_B+1$ accounts precisely for the guard edge.

\paragraph{Sublevel boundedness.}
Fix $r\geq0$ and choose $H_r\in\mathbb N$ bounding $u$ on
$\{\sigma\in\supp(I)\mid v(\sigma)\leq r\}$.
At guard configurations with $V\leq r$, we have $U\leq aH_r+1$;
at the program exit, $U=0$. At an internal configuration $\gamma$,
backward expansion expresses $V(\gamma)$ as the expectation of
$v(\sigma')+1$ at the next guard. This is a finite probability
distribution. If $V(\gamma)\leq r$, at least one positive-probability
endpoint satisfies $v(\sigma')+1\leq r$. Choose a path to that endpoint.
The path inequality above yields
\[
 U(\gamma)\leq a\,u(\sigma')+1+L_B\leq aH_r+1+L_B.
\]
This is a uniform bound over the entire pCFG sublevel set, including
internal locations. All premises of \Cref{thm:pcfg-rule} now hold,
with the same progress constant $p_*$ for every sublevel set. Its
soundness and AST preservation by the translation establish AST of
$C$ on $M$. This argument is pointwise in initial states; it requires
no finite expectation of $v$ under an arbitrary $\mu\in M$.
\end{proof}

\subsection{Relative Completeness}

\paragraph{Assertion language.}
For relative completeness, we use first-order arithmetic with addition,
multiplication, and order interpreted over $\mathbb Q$, as
in~\cite{DBLP:journals/pacmpl/MajumdarS25}. Program states have effective
rational encodings, and guards and updates are total computable operations
on these encodings. Initial supports are arithmetically definable.
Natural numbers and computable relations are definable in this rational
arithmetic~\cite[Section~2.4]{DBLP:journals/pacmpl/MajumdarS25}.
Natural-valued certificates are represented by their graphs, and
real-valued certificates by predicates for rational cuts
$q<f(\sigma)$; they need not be rational-valued or arithmetic terms.
For distributional invariants we allow the lifted support assertion
$\supp(\mu)\subseteq S$ for an arithmetical state predicate $S$.
Completeness is relative to all true sentences of this arithmetic
structure, not to a decision procedure or to the templates of
\Cref{sec:certificate-synthesis}.

\begin{theorem}[Relative completeness]
\label{thm:complete}
Let $C=\while{b}{B}$ have loop-free body. If $C$ is AST on $M$, then
certificates satisfying \Cref{thm:program-rule} exist. Under the preceding
arithmetic assumptions, they are expressible in the assertion language,
and the rule is relatively complete.
\end{theorem}
\begin{proof}
Put $E=\supp(M)$. By AST preservation, $\mathcal T[C]$ is AST from
every state in $E$. The pCFG completeness result supplies common
certificates $\mathit{Inv},V,U$ for these initial states; the extension
from one initial state to a set of initial states is justified in
\Cref{app:initial-sets}. Set
\begin{equation}
 \begin{split}
 S&=\{\sigma\mid(\ell_g,\sigma)\in\mathit{Inv}\},\qquad
 I=\{\mu\in\subdists\mid\supp(\mu)\subseteq S\},\\
 v(\sigma)&=[b](\sigma)V(\ell_g,\sigma),\qquad
 u(\sigma)=[b](\sigma)U(\ell_g,\sigma)\quad(\sigma\in S).
 \end{split}
 \label{eq:project-certificates}
\end{equation}
Inductiveness of $\mathit{Inv}$ and the operational correspondence show
that $S$ contains $E$ and is closed under enabled body executions.
Support closure therefore makes $I$ a distributional invariant, with
$\supp(I)=S$.

The projected functions have the required codomains and vanish on
$S\cap\llbracket\neg b\rrbracket$; on enabled states $v$ is positive.
For such a state $\sigma$, \Cref{lem:contract-v} and
$v(\sigma')\leq V(\ell_g,\sigma')$ give
\[
 \sum_{\sigma'}\post{B}(\dirac{\sigma})(\sigma')v(\sigma')
 \leq\sum_{\sigma'}\post{B}(\dirac{\sigma})(\sigma')
                         V(\ell_g,\sigma')
 \leq V(\ell_g,\sigma)=v(\sigma).
\]
Similarly, every endpoint counted in \eqref{eq:contract-u} satisfies
\[
 u(\sigma')\leq U(\ell_g,\sigma')
   <U(\ell_g,\sigma)=u(\sigma).
\]
Thus the source decrease premise holds with the single constant
$\varepsilon_B$. Finally, for every $r\geq0$,
\[
 \{u(\sigma)\mid\sigma\in S,\ v(\sigma)\leq r\}
 \subseteq\{0\}\cup
 \{U(\gamma)\mid\gamma\in\mathit{Inv},\ V(\gamma)\leq r\}.
\]
The right-hand side is bounded. This inclusion handles terminal states
without incorrectly identifying their $v$- and $V$-sublevel membership.

For expressibility, fixing the location variable in the formula for
$\mathit{Inv}$ defines $S$, and hence its lifted support assertion.
The graph of $u$ is obtained by a guard case distinction. If
$\theta_V(\ell,\sigma,q)$ defines $q<V(\ell,\sigma)$, the cut of $v$
is defined on $S$ by
\[
 \bigl(b(\sigma)\land\theta_V(\ell_g,\sigma,q)\bigr)
 \lor\bigl(\neg b(\sigma)\land q<0\bigr).
\]
All premises reduce to arithmetic sentences: each body has finitely many
symbolic paths with rational probabilities and computable updates, so
its sums are finite; comparisons of real values use their rational cuts.
For sublevel boundedness it suffices to quantify over integer thresholds
$r$ and integer bounds. Invariant initialization and preservation reduce
to $E\subseteq S$ and the support-closure condition. Their established
validity therefore supplies a proof relative to true arithmetic.
\end{proof}

The two constructions preserve semantic proving power for this fragment.
Users supply only the two source-state functions; the internal annotations
are needed solely to justify the metatheory.
